\documentclass[10pt,a4paper]{amsart}
\usepackage[margin=19mm]{geometry}
\usepackage{amsmath,amssymb,mathtools}
\usepackage[scr=rsfs,cal=euler]{mathalfa}
\usepackage{xfrac}
\usepackage[unicode,hidelinks,bookmarks=false]{hyperref}
\newcommand{\ie}{i.e.\ }
\newcommand{\cf}{cf.\ }
\newcommand{\resp}{respectively\ }
\newcommand{\Subsection}{\subsection}
\providecommand{\nobreakdash}{\nobreak\hbox{-}}
\setlength{\emergencystretch}{2em}
\multlinegap0pt
\usepackage{bm}
\usepackage{bbm}
\usepackage{dsfont}
\usepackage{xspace}
\usepackage{cancel}
\usepackage{mathtools}
\usepackage{amsthm}
\makeatletter
\@ifundefined{theorem}{\newtheorem{theorem}{Theorem}}{}
\@ifundefined{lemma}{\newtheorem{lemma}[theorem]{Lemma}}{}
\makeatother
\title[Longest paths in trees and isometricity of ultrametric space]{Longest paths in trees and isometricity of ultrametric spaces}
\author{Oleksiy Dovgoshey, Olga Rovenska}
\date{Source: arXiv:2510.10038v1 (CC BY 4.0)}
\begin{document}
\maketitle

\noindent\emph{Source and licence.} Longest paths in trees and isometricity of ultrametric spaces. Version arXiv:2510.10038v1 (CC BY 4.0), distributed under the Creative Commons Attribution 4.0 International licence, as recorded in the arXiv licence metadata for this exact version and shown on the abstract page https://arxiv.org/abs/2510.10038v1. Full rights notice: licenses/arxiv-2510.10038-CC-BY-4.0.txt.\par\smallskip

\noindent\textit{Editorial scope.} Counted unit: the Lemma whose source label is \texttt{wsdt} in the article (source0.tex lines 464--482, source label \texttt{wsdt}), delivered together with the article's own complete original proof of it (source0.tex lines 484--535, one \texttt{proof} environment of 52 lines). No mathematics is written, repaired or re-derived: statement and proof are the article's own text line for line. Required context retained uncounted: none. Every definition, display and label of those lines is retained unchanged. Omitted from this package and not redistributed: 1--463, 483--483, 536--971. Nothing inside a retained excerpt is omitted or altered: every retained body is delivered exactly as the article's lines are written, so no omission receipt of any kind accompanies this package. Citations: the retained excerpts carry no citation command at all, so no bibliography entry of the article is transcribed and no reference is redistributed. Reference adaptation: none was needed and none was made; every reference inside the delivered unit resolves inside it, so no unresolved reference or citation is produced. Printed slips kept literal: no printed word, symbol, hypothesis or number of the counted statement or of its proof was altered. Layout and class: the article's own class is \texttt{article} with packages including enumitem, lineno, amsfonts, amsmath, amsthm, amscd, amssymb, latexsym, amsbsy, pb-diagram, cite, caption, url, hyperref; the wrapper is the lane's pinned \texttt{amsart} editorial wrapper at 10pt on A4, which restates the theorem-like environments the retained text needs and adds the article's own macro definitions that the retained text needs (none), with extra wrapper packages (bm, bbm, dsfont, xspace, cancel, mathtools). The delivered numbering is the unit's own. Assets: no retained span contains a figure environment, a graphics-inclusion command, a PDF-page inclusion, a table or a TikZ picture; no image file is copied into this package, the package declares no asset dependency and no image file is redistributed with it. Licence: version arXiv:2510.10038v1 of this article is distributed under the Creative Commons Attribution 4.0 International licence, as recorded in the arXiv licence metadata for this exact version and shown on the abstract page https://arxiv.org/abs/2510.10038v1. A full rights notice is delivered at licenses/arxiv-2510.10038-CC-BY-4.0.txt. Characters: every retained excerpt is pure ASCII, so the delivered engine, XeLaTeX with its default Latin Modern text font, reports no missing character.
\begin{lemma}\label{wsdt}
Let \( T \) be a tree and let 
\begin{equation}
    \label{rfhi}
    |E(P)| \leq 3
\end{equation}
holds for every path $P\subseteq T$.  
Then 
every two distinct vertices \( u\) and \( v \) whose degree is greater than one, 
\begin{equation}
    \label{hhij}
\operatorname{deg} u \geq 2 \quad \text{and} \quad \operatorname{deg} v \geq 2,
\end{equation}
are adjacent,  
\begin{equation}
    \label{tyhj}
\{u,v\} \in E(T). 
\end{equation}
\end{lemma}

\begin{proof}
Let distinct $u,v\in V(T)$ satisfy \eqref{hhij}.
Suppose that
\begin{equation}
    \label{rvg}
    \{u,v\} \notin E(T). 
\end{equation}
Since every tree is a connected graph, relation~\eqref{rvg} implies the existence of a path \( P^0 \subseteq T \) such that
$
P^0= (v_1, \ldots, v_n),
$
with \( v_1 =u,\) \( v_n = v \) and \( n \geq 3 \).
Using inequalities \eqref{hhij} we can find some vertices $v_0, v_{n+1}$ of $T$ such that
\begin{equation}
    \label{ds1}
\{v_0, v_1\} \in E(T), \quad 
\{v_{n}, v_{n+1}\} \in E(T)
\end{equation}
and 
\begin{equation}
     \label{ds2}
     v_0 \neq v_2, \quad v_{n+1} \neq v_{n-1}.
\end{equation}
The tree $T$ is an acyclic graph. Consequently we also have
\begin{equation}     \label{ds3}
v_0 \notin \{v_2, \ldots, v_n\}
\quad {\rm and} \quad
     v_{n+1} \notin \{v_1, \ldots, v_{n-2}\}.
\end{equation}
Now relations \eqref{ds1}--\eqref{ds3} imply that
$
P^* = (v_0, v_1, \ldots, v_n, v_{n+1})
$
is a path joining $v_0$ and $v_{n+1}$ in $T$.
It is clear that $P^0\subseteq P^*\subseteq T$ and
\begin{equation}
    \label{ds4}
|E(P^*)| = |E(P^0)| + 2. 
\end{equation}
Furthermore, the inequality
\begin{equation*}
    |E(P^0)| \geq 2
\end{equation*}
holds by the definition of the path $P^0$. The last inequality and \eqref{ds4} give us the inequality
\[
|E(P^*)| \geq 4,
\]
that contradicts \eqref{rfhi}. 

Relation \eqref{tyhj} follows. 

\end{proof}

\end{document}
